% Architecture fonctionnelle globale d'AutoMaze AI (cinq étages, de l'environnement au contrôleur)
\begin{tikzpicture}[
    etage/.style={draw=accent, fill=accent!7, rounded corners=3pt, align=center, font=\footnotesize,
                  text width=12.2cm, minimum height=1.0cm, inner sep=4pt},
    bloc/.style={draw=black!60, fill=white, rounded corners=3pt, align=center, font=\footnotesize,
                 text width=3.6cm, minimum height=1.05cm, inner sep=3pt},
    ml/.style={bloc, draw=accent, fill=accent!18, very thick},
    fleche/.style={-{Stealth[length=2.2mm]}, draw=black!70, thick}
  ]
  \node[etage] (env) at (0,0) {\textbf{1. Environnement de simulation}\\Routes · Véhicule · Objets routiers · Destination};

  \node[bloc] (camera) at (-4.3,-1.75) {\textbf{Caméra virtuelle}\\Image de l'objet};
  \node[ml] (modele) at (-4.3,-3.35) {\textbf{Modèle ML}\\Classification de l'objet};
  \node[bloc] (infos) at (4.3,-2.55) {\textbf{Informations du simulateur}\\Coordonnées et positions des objets};
  \node[bloc] (etat) at (0,-4.75) {\textbf{État des obstacles}\\Position + nature identifiée};

  \node[etage] (fusion) at (0,-6.3) {\textbf{2. Fusion des informations et représentation de l'environnement}\\Association de la classe prédite aux coordonnées de l'objet correspondant};
  \node[etage] (plan) at (0,-7.75) {\textbf{3. Planificateur}\\Graphe · Dijkstra · Recalcul du trajet};
  \node[etage] (decision) at (0,-9.2) {\textbf{4. Système de décision}\\Ralentir · Arrêter · Contourner · Continuer};
  \node[etage] (controle) at (0,-10.65) {\textbf{5. Contrôleur du véhicule}\\Vitesse · Freinage · Direction};
  \node[font=\footnotesize\bfseries, text=integration] (but) at (0,-11.75) {Déplacement autonome jusqu'à la destination};

  \draw[fleche] (env.south -| camera) -- (camera);
  \draw[fleche] (camera) -- (modele);
  \draw[fleche] (env.south -| infos) -- (infos);
  \draw[fleche] (modele.south) |- (etat.west);
  \draw[fleche] (infos.south) |- (etat.east);
  \draw[fleche] (etat) -- (fusion);
  \draw[fleche] (fusion) -- (plan);
  \draw[fleche] (plan) -- (decision);
  \draw[fleche] (decision) -- (controle);
  \draw[fleche] (controle) -- (but);

  % Une décision qui bloque un tronçon relance le planificateur
  \draw[fleche, dashed, draw=accent] (decision.east) -- ++(0.55,0) |- node[pos=0.25, right, font=\scriptsize, align=left, text=accent] {obstacle\\bloquant :\\recalcul} (plan.east);
  % Le véhicule se déplace dans l'environnement : la boucle reprend
  \draw[fleche, draw=integration] (but.west) -- ++(-6.9,0) |- node[pos=0.25, left, font=\scriptsize, text=integration, align=right] {nouvelle\\position} (env.west);

  % Ce que le simulateur ne transmet pas
  \node[font=\scriptsize, text=black!65, align=center, text width=3.8cm] at (4.3,-3.95) {La vraie classe de l'objet n'est jamais transmise à la décision};
\end{tikzpicture}
